\section{Forward/backward kinetics와 인과적 기억}\label{KIN-CH3}

평형식은 target을 정하지만 그 target에 도달하는 시간을 정하지
않는다. 이 장은 Chapter~\ref{EQ-CH1}의 상태축을 유지한 채
forward/backward flux, local detailed balance, 전이상태 속도와
인과적 relaxation을 정의한다.

\subsection{두 방향 flux와 stationary target}\label{KIN-CH3-RATES}

\paragraph{\physid{KIN-001} mass-action skeleton.}
\begin{equation}
\boxed{
\dot\xi_j=J_j^+-J_j^-,
\qquad
J_j^+=r_j^+(1-\xi_j),
\qquad
J_j^-=r_j^-\xi_j.
}
\label{KIN-001-EQUATION}
\end{equation}
$r_j^\pm\ge0$의 단위는 s$^{-1}$다. 방향은
$J^+-J^-$가 담당하며 branch마다 $\xi_j$의 뜻을 바꾸지 않는다.

\paragraph{\physid{KIN-002} explicit target의 조건.}
$r_j^\pm$가 $\xi_j$에 무관하거나 한 local step에서 고정될 때만
\begin{equation}
\xi_{\rm ss,j}=\frac{r_j^+}{r_j^++r_j^-},
\qquad
k_j=r_j^++r_j^-
\label{KIN-002-EXPLICIT}
\end{equation}
이고
\[
\dot\xi_j=k_j(\xi_{\rm ss,j}-\xi_j)
\]
로 정확히 재배열된다. $r_j^\pm=r_j^\pm(\xi_j)$이면 stationary
state는
\begin{equation}
r_j^+(\xi_{\rm ss})(1-\xi_{\rm ss})
=r_j^-(\xi_{\rm ss})\xi_{\rm ss}
\label{KIN-002-IMPLICIT}
\end{equation}
의 implicit root다. 식~\eqref{KIN-002-EXPLICIT}를 그대로
대입하지 않는다.

\subsection{local detailed balance와 화학양론}\label{KIN-CH3-BALANCE}

\paragraph{\physid{KIN-003} flux ratio와 affinity.}
같은 coarse-grained edge의 두 unidirectional flux가 양수이면
\begin{equation}
\boxed{
\mathcal A_j^\chem
=RT\ln\frac{J_j^+}{J_j^-}.
}
\label{KIN-003-EQUATION}
\end{equation}
평형에서는 $J_j^+=J_j^-$이고 $\mathcal A_j^\chem=0$이다.
식~\eqref{KIN-003-EQUATION}는 Chapter~\ref{THM-CH2}의
entropy production~\eqref{THM-006-NETWORK}과 같은 공액쌍을
사용한다.

\paragraph{\physid{KIN-004} center coordinate와 affinity의 구분.}
$z_jF(V_{\drive}-U_j)$는 ideal two-state closure에서 affinity의
전위 부분이 될 수 있으나, 일반 상태에서는 configurational,
stress와 activity 기여가 추가된다. 따라서
\[
z_jF(V_{\drive}-U_j)
\]
를 모든 상태의 chemical affinity로 선언하지 않는다.
$U_j$ 중심 기준 좌표는 mobility를 정리하는 보조좌표이고,
진짜 affinity는 평형에서 0이 되는 식~\eqref{KIN-003-EQUATION}다.

\paragraph{\physid{KIN-005} $z_j$와 폭의 독립.}
\[
\boxed{z_j\not\equiv\frac{RT}{Fw_j}.}
\]
$z_j$는 balanced reaction이 정하며 Li 한 개의 삽입/탈리 반응은
보통 $z_j=1$이다. 이상 독립-site closure에서만
$w_{\rm th}=RT/(z_jF)$가 따라온다. 정적 broadening이나 경험적
$w_j$를 이 식에 넣어 전자수를 재정의하지 않는다.

\subsection{전이상태 속도와 구동력}\label{KIN-CH3-EYRING}

\paragraph{\physid{KIN-006} Eyring prefactor.}
\begin{equation}
k_0(T)=\frac{k_BT}{h}\kappa(T),\qquad
k^\pm=k_0^\pm
\exp\!\left[-\frac{\Delta G^{\ddagger,\pm}}{RT}\right].
\label{KIN-006-EQUATION}
\end{equation}
$\kappa(T)$는 transmission coefficient다. 독립적으로 알지 못하면
activation entropy와 prefactor intercept에서 분리되지 않는다.
$k_BT/h$만을 보편적인 완전 prefactor로 두지 않는다.

\paragraph{\physid{KIN-007} 방향성 barrier split.}
작은 affinity에서 한 가능한 local closure는
\begin{align}
\Delta G^{\ddagger,+}
&=\Delta G_0^{\ddagger,+}-\beta\mathcal A^\chem,\\
\Delta G^{\ddagger,-}
&=\Delta G_0^{\ddagger,-}+(1-\beta)\mathcal A^\chem,
\qquad 0\le\beta\le1 .
\end{align}
이는 small-driving-force 1차 근사다. rate ratio는 target을,
rate sum은 mobility를 정한다. 공통 mobility sensitivity와
directional split $\beta$를 동일 파라미터로 취급하지 않는다.
강구동에서는 Marcus 곡률, 수송, site availability와 current
allocation을 포함한 별도 closure가 필요하다.

\paragraph{\physid{KIN-008} 물리적 병목.}
독립적인 직렬 단계가 실제로 존재할 때만 시간척도의 합
\[
\frac1{k_{\rm phys}}
=\frac1{k_{\rm act}}+\frac1{k_{\rm tr}}+\frac1{k_{\rm site}}
\]
을 쓴다. 이것은 임의의 rate cap이 아니다. 한 dataset에서
$k_{\rm act}$와 $k_{\rm tr}$를 모두 자유화하면 강하게 교락되므로,
각 항에는 독립적인 rate dependence 또는 transient evidence가
필요하다.

\subsection{SI 시간계약}\label{KIN-CH3-UNITS}

\paragraph{\physid{KIN-009} 초 단위 속도.}
물리 속도식의 $r^\pm,k,\dot q$는 s$^{-1}$로 통일한다. C-rate
$C_h$가 h$^{-1}$이면
\[
\dot q=\frac{C_h}{3600}\quad[\mathrm{s^{-1}}].
\]
3600은 온도에 무관한 rate multiplier다. 따라서 다온도
$\ln(k/T)$ 대 $1/T$의 slope는 단위변환 자체로 바뀌지 않고
intercept가 바뀐다. 고정 $\Delta H_a$ convention에서는
\begin{equation}
\Delta S_a^{\rm SI}-\Delta S_a^{\rm hour}
=-R\ln3600=-68.0808\ {\rm J\,mol^{-1}K^{-1}}.
\label{KIN-009-EQUATION}
\end{equation}
단일온도에서 $\Delta S_a$를 고정해 offset을 $\Delta H_a$에
흡수한 경우의 $RT\ln3600$은 그 온도에서의 apparent correction일
뿐, 다온도 physical enthalpy correction이 아니다.

\subsection{인과적 relaxation과 prehistory}\label{KIN-CH3-MEMORY}

\paragraph{\physid{KIN-010} 시간 trajectory.}
\begin{equation}
\dot\xi_j=\frac{\xi_{\tar,j}(t)-\xi_j}{\tau_j(t)}
\label{KIN-010-EQUATION}
\end{equation}
는 acquisition-time 순서에서 적분한다. reversal, rest 또는
비단조 전압을 전압순으로 정렬하면 prehistory를 잃으므로
식~\eqref{KIN-010-EQUATION}과 다른 문제가 된다.

\paragraph{\physid{KIN-011} 초기조건 계약.}
상수 $\tau$의 해는
\begin{equation}
\xi(t)=\xi(t_0)e^{-(t-t_0)/\tau}
+\int_{t_0}^{t}\frac{e^{-(t-t')/\tau}}{\tau}
\xi_{\tar}(t')\,\dd t'.
\label{KIN-011-EQUATION}
\end{equation}
따라서 다음 중 하나가 필요하다: 측정된 prehistory, supplied
initial state, 명시적 asymptotic prehistory 또는 오차가 표시된
finite-padding approximation. $5\tau$ padding에는 여전히
$e^{-5}\simeq0.00674$의 잔여가 있으므로 $-\infty$ 경계와 같다고
쓰지 않는다.

\paragraph{\physid{KIN-012} monotonic curve 좌표.}
정전류이며 한 방향으로 단조인 구간에서만
\[
\frac{\dd\xi}{\dd q}
=\frac{Q_\cell}{|I|}
\frac{\xi_{\tar}-\xi}{\tau}
\]
로 바꿀 수 있다. 이때 $L_q=|I|\tau/Q_\cell$는 무차원
relaxation length다. 휴지의 $I=0$ 또는 reversal을 이 좌표식으로
통과시키지 않는다.

\paragraph{\physid{KIN-013} spectrum의 정규화와 진폭.}
분포 relaxation은
\[
\int_0^\infty a(L)\,\dd L=1,\qquad
\Theta(q)=\int_0^\infty a(L)\xi_L(q)\,\dd L
\]
처럼 normalized measure $a(L)$를 먼저 둔다. tail 시작점의
mode-dependent residual amplitude는 별도 $b(L)$로 둔다.
$a(L)$와 $b(L)a(L)$를 같은 기호로 써서 probability와 관측
진폭을 혼동하지 않는다.

\subsection{식별성과 전달}\label{KIN-CH3-IDENT}

\begin{identifiability}
단일온도ㆍ단일 C-rate curve에서는 $\kappa$, $\Delta H^\ddagger$,
$\Delta S^\ddagger$, $\beta$, transport limitation, static
broadening과 relaxation spectrum을 고유하게 분리할 수 없다.
다온도 slope, rate series, current-interrupt/GITT transient와
독립적인 수송자료가 필요하다.
\end{identifiability}

이 장은 $\dot{\bm\xi}=\bm f(\bm\xi,V_n,T,I)$와 그 initial-history
계약을 Chapter~\ref{BAL-CH4}에 넘긴다. constitutive closure가
아직 없는 항은 열린 모듈로 남는다.
